% Els criteris generals són a la pàgina 1 de la pauta, sota l'encapçalament de la
% sèrie 3; la part de la sèrie 5 (pàgines 9-13) no en té de propis, i s'entén que
% són comuns a totes les sèries.
% Els criteris 5, 6, 7 i part del 8 apareixen subratllats a l'original.
\begin{enumerate}
  \item Les respostes s'han d'ajustar a l'enunciat de la pregunta. Es valorarà sobretot que l'alumnat demostri que té clars els conceptes de caràcter físic sobre els quals tracta cada pregunta.
  \item Es tindrà en compte la claredat en l'exposició dels conceptes, dels processos, dels passos a seguir, de les hipòtesis, l'ordre lògic, l'ús correcte dels termes científics i la contextualització segons l'enunciat.
  \item En les respostes cal que l'alumnat mostri una adequada capacitat de comprensió de les qüestions plantejades i organitzi de forma lògica la resposta, tot analitzant i utilitzant les variables en joc. També es valorarà el grau de pertinença de la resposta, el que l'alumnat diu i les mancances manifestes sobre el tema en qüestió.
  \item Totes les respostes s'han de raonar i justificar. Una resposta correcta sense raonament ni justificació pot ser valorada amb un zero si el corrector no és capaç de veure d'on ha sortit el resultat, tanmateix un resultat erroni amb un raonament correcte s'ha de valorar.
  \item Un error no s'ha de penalitzar dues vegades en el mateix problema. Si un apartat necessita un resultat anterior, i aquest és erroni, cal valorar la resposta independentment del seu valor numèric, i tenir en compte el procediment de resolució.
  \item Si l'alumne ha resolt un problema per un altre procediment vàlid diferent del descrit en aquestes pautes, la resolució es considera vàlida.
  \item Els errors d'unitats o el fet de no posar-les restaran el 20\,\% de la puntuació d'aquest apartat. Exemple: Si un apartat val 1 punt i s'ha equivocat en les unitats li haurem de puntuar 0,8 punts.
  \item Cal resoldre els exercicis fins al resultat final i no es poden deixar indicades les operacions. Tanmateix, els errors en el càlcul restaran el 20\,\% de la puntuació d'aquest apartat. Exemple: Si un apartat val 1 punt i s'ha equivocat en els càlculs li haurem de puntuar 0,8 punts.
  \item Cal fer la substitució numèrica a les expressions que s'usen per resoldre les preguntes.
  \item Un resultat amb un nombre molt elevat de xifres significatives (6 xifres significatives) es penalitzarà amb 0,1\,p.
\end{enumerate}
